\subsection{Other Selection Criteria}\label{sec:entropy-selection-criteria}

The Lax shock condition is not the only viable selection criterion used to pick the "physically reasonable" solution from among the possible weak solutions to systems of conservation laws. In this section we present some competing conditions and describe some of the relationships between them.

\subsubsection{The Entropy Condition}\label{subsec:entropy-condition}

The first alternative selection criterion we present is called the entropy condition. It is an outgrowth of the second law of thermodynamics, which is generalized in this situation to include physical systems other than mechanical and thermal. The key to the condition is the existence of an additional conservation law derived from \eqref{eq:conservation-law-system}.

\begin{Definition}
\label{def:entropy-flux-pair}
An \textbf{entropy/entropy-flux pair}\footnote{This terminology is quite standard despite the fact that it contradicts any reasonable definitions from thermodynamics. Perhaps a better term would be free energy/energy-flux, but it seems hopeless to try to change things after so many years. Students of conservation laws should therefore resign themselves to forgetting everything they have learned about thermodynamics while studying the subject. The culprit responsible for this terminology has not come forward to confess.} is a pair of functions $(U,F):\mathbb{R}^n\to\mathbb{R}^2$ satisfying
\begin{equation}\label{eq:entropy-flux-pair-definition}
\nabla F=\nabla U\cdot\nabla\bm{f}.
\end{equation}

\end{Definition}

It follows immediately from the definition and the chain rule that if $\bm{u}$ is a classical solution of \eqref{eq:conservation-law-system}, then
\begin{equation}\label{eq:classical-entropy-conservation}
U(\bm{u})_t+F(\bm{u})_x=0.
\end{equation}
Of course, as we noted in Remark~\ref{rem:weak-solution-variable-change}, a weak solution of \eqref{eq:conservation-law-system} does not necessarily satisfy \eqref{eq:classical-entropy-conservation}.

\begin{Definition}
\label{def:entropy-condition}
A weak solution of \eqref{eq:conservation-law-system} is said to satisfy the \textbf{entropy condition} if there exists an entropy/entropy-flux pair with $\bm{u}\mapsto U(\bm{u})$ convex such that
\begin{equation}\label{eq:entropy-weak-inequality}
-\int\!\!\int(U(\bm{u})\phi_t+F(\bm{u})\phi_x)\,\mathrm{d}x\,\mathrm{d}t\leq0
\end{equation}
for every \emph{non-negative} $C^1$ test function $\phi$ with compact support in the open upper half-plane.
\end{Definition}

\begin{Remark}
\label{rem:entropy-distribution-interpretation}
An entropy/entropy-flux pair satisfying \eqref{eq:entropy-weak-inequality} is often described as satisfying the inequality
\begin{equation}\label{eq:entropy-distribution-inequality}
U(\bm{u})_t+F(\bm{u})_x\leq0
\end{equation}
in the \emph{sense of distributions}. Note that if an entropy/entropy-flux pair satisfies \eqref{eq:entropy-distribution-inequality} in the classical sense, then we can multiply the inequality by a non-negative test function and derive \eqref{eq:entropy-weak-inequality} using Green\textquotesingle s theorem in the plane.
\end{Remark}

\begin{Remark}
\label{rem:entropy-free-energy}
If $\bm{u}$ satisfies \eqref{eq:entropy-distribution-inequality} and has compact support, then by integrating \eqref{eq:entropy-distribution-inequality} with respect to $x$ we get
\begin{equation}\label{eq:total-entropy-time-derivative}
\int_{-\infty}^\infty U_t\,\mathrm{d}x\leq0.
\end{equation}
Thus, the quantity
\[
\int_{-\infty}^\infty U\,\mathrm{d}x
\]
(which, we repeat, is best thought of as a free energy, rather than an entropy) is a decreasing function of time. The idea that energy can only be lost can be thought of as a version of the second law of thermodynamics.
\end{Remark}

The following lemma translates the entropy condition into a jump condition for piecewise continuous weak solutions.

\begin{Lemma}
\label{lem:entropy-jump-condition}
Suppose $\bm{u}$ is a piecewise continuous weak solution of \eqref{eq:conservation-law-system} that satisfies the entropy condition. Suppose $\bm{u}$ has a jump discontinuity along a shock curve with slope $s$. Then
\begin{equation}\label{eq:entropy-jump-condition}
s[U(\bm{u})]-[F(\bm{u})]\geq0.
\end{equation}

\end{Lemma}

The proof of this is left to the reader (cf. Problem~\ref{ex:entropy-jump-condition-proof}). The proof can be obtained by modifying the proof of the Rankine-Hugoniot condition.

The following theorem states that in the case of genuinely nonlinear systems, the entropy condition and the Lax shock condition are equivalent for weak shocks.

\begin{Theorem}
\label{thm:entropy-lax-small-shock-equivalence}
Suppose that the system of conservation laws \eqref{eq:conservation-law-system} is strictly hyperbolic and genuinely nonlinear. In addition, suppose that there is an entropy/entropy-flux pair $(U,F)$ with $U$ strongly convex. Let $\bm{u}$ be a weak solution of \eqref{eq:conservation-law-system}, containing a shock with values $\bm{u}^l$ and $\bm{u}^r$ traveling with speed $s$, with $\bm{u}^l-\bm{u}^r$ sufficiently small. Then $\bm{u}$ satisfies the Lax shock condition at this shock if and only if the entropy jump condition \eqref{eq:entropy-jump-condition} holds with strict inequality.
\end{Theorem}

\begin{Proof}
Since $\bm{u}$ is a weak solution, $\bm{u}^l$, $\bm{u}^r$ and $s$ must satisfy the Rankine-Hugoniot condition. Thus, by Theorem~\ref{thm:shock-curves-existence}, $\bm{u}^r$ must lie along one of the $n$ shock curves through $\bm{u}^l$; i.e., $\bm{u}^r=\hat{\bm{u}}_k(\epsilon)$ and $s=\hat s_k(\epsilon)$ for some $k=1,2,\ldots,n$ and some $\epsilon\in I_k$. Now, again by Theorem~\ref{thm:shock-curves-existence}, the Lax shock condition is satisfied if and only if $\epsilon<0$. Furthermore, the entropy jump condition holds if and only if
\begin{equation}\label{eq:entropy-shock-curve-functional}
E(\epsilon):=\hat s_k(\epsilon)[U(\bm{u}^l)-U(\hat{\bm{u}}_k(\epsilon))]-[F(\bm{u}^l)-F(\hat{\bm{u}}_k(\epsilon))]<0.
\end{equation}
We now let a \textquotedblleft prime\textquotedblright\ indicate differentiation with respect to $\epsilon$ and show $E(0)=E'(0)=E''(0)=0$ and that $E'''(0)>0$. This indicates that for $\epsilon$ sufficiently small, $E$ is negative if and only if $\epsilon$ is negative and thus completes the proof. Using $\bm{u}^l=\hat{\bm{u}}_k(0)$ we get
\begin{equation}\label{eq:entropy-shock-functional-at-zero}
E(0)=\hat s_k(0)[U(\bm{u}^l)-U(\hat{\bm{u}}_k(0))]-[F(\bm{u}^l)-F(\hat{\bm{u}}_k(0))]=0.
\end{equation}
We could use $\hat s_k(0)=\lambda_k$ and $\bm{u}_k'(0)=\bm{r}_k$ to calculate $E'(0)$ directly, but instead we differentiate the Rankine-Hugoniot condition to get
\begin{equation}\label{eq:shock-rh-parameter-derivative}
\hat s_k'(\epsilon)[\bm{u}^l-\hat{\bm{u}}_k(\epsilon)]-\hat s_k(\epsilon)\hat{\bm{u}}_k'(\epsilon)=-\bm{f}(\hat{\bm{u}}_k(\epsilon))'.
\end{equation}
We now use $\nabla F=\nabla U\cdot\nabla\bm{f}$ to get
\begin{equation}\label{eq:entropy-shock-functional-first-derivative}
E'(\epsilon)=\hat s_k'(\epsilon)\{[U(\bm{u}^l)-U(\hat{\bm{u}}_k(\epsilon))]-\nabla U(\hat{\bm{u}}_k(\epsilon))\cdot[\bm{u}^l-\hat{\bm{u}}_k(\epsilon)]\},
\end{equation}
from which it is easy to see that $E'(0)=0$. Simply differentiating this gives us $E''(0)=0$. The calculation of $E'''(0)$ contains many terms that go to zero in the same way as the terms of the preceding calculations, but one interesting term remains:
\begin{equation}\label{eq:entropy-shock-functional-third-derivative}
E'''(0)=\hat s_k'(0)[\hat{\bm{u}}_k'(0)^T\nabla^2U(\hat{\bm{u}}_k(0))\hat{\bm{u}}_k'(0)]
\end{equation}
where $\nabla^2U$ is the second gradient or Hessian matrix of $U$. Now, from Theorem~\ref{thm:shock-curves-existence} we have $\hat s_k'(0)=1/2$ and since $U$ is strictly convex its Hessian is positive definite. Thus $E'''>0$, and the theorem is proved.
\end{Proof}
